% ============================================================
%  Regra de Três e Porcentagem — Apostila Premium | PratiqueJá
%  Engine: XeLaTeX
% ============================================================
\documentclass[12pt]{article}
\usepackage{../../pratiqueja}

% \hlm — destaque de resultado em modo-math (cor sem bold)
\newcommand{\hlm}[1]{{\color{darkblue}#1}}

\setsubject{Regra de Três e Porcentagem}

\begin{document}
\pagenumbering{arabic}
\setstretch{1.2}
\color{bodytext}

\heroheader{Regra de Três e Porcentagem}%
           {Proporções Simples, Compostas e \%}%
           {Matemática · Básico}

\setlength{\columnsep}{18pt}
\setlength{\columnseprule}{0.5pt}
\def\columnseprulecolor{\color{exborder}}

\begin{multicols}{2}

%─────────────────────────────────────────────────────────────
\TW{Def}{badgepurple}{Definição}{Razão, proporção e regra de três}

\regra{Uma \textbf{proporção} é a igualdade entre duas razões:
$\dfrac{a}{b} = \dfrac{c}{d}$. A regra de três acha um valor
desconhecido $x$ a partir de três conhecidos.}

\small\color{bodytext}%
Com \textbf{duas} grandezas: regra de três \textbf{simples}. Com
\textbf{três ou mais}: \textbf{composta}.

\nota{\textbf{Como identificar o tipo:} ``Se A aumenta, B aumenta ou
diminui?''\\
Aumenta → \textbf{direta} $(\downarrow\downarrow)$\\
Diminui → \textbf{inversa} $(\downarrow\uparrow)$}

%─────────────────────────────────────────────────────────────
\TW{$\times$}{darkblue}{Produto Cruzado}{Método geral de resolução}

\regra{Em $\dfrac{a}{b} = \dfrac{c}{d}$, o \textbf{produto dos
extremos é igual ao dos meios}:}
\formula{$a \cdot d = b \cdot c$}

\exemp{%
  $\dfrac{2}{x} = \dfrac{1}{3} \;\Rightarrow\; 2\cdot3 = 1\cdot x
   \;\Rightarrow\; x = \hlm{6}$\\[4pt]
  $\dfrac{x}{6} = \dfrac{2}{4} \;\Rightarrow\; 4x = 2\cdot6
   \;\Rightarrow\; x = \dfrac{12}{4} = \hlm{3}$%
}

%─────────────────────────────────────────────────────────────
\TW{Dir}{darkblue}{Simples — Proporção Direta}{Grandezas que crescem na mesma razão}

\regra{\textbf{Diretamente proporcionais:} o aumento de uma provoca
o aumento da outra na mesma razão.}

\SH{Exemplo 1 — horas e peças}
1 peça em 2 horas; quantas horas para 3 peças? \emph{Mais peças,
mais horas → direta.}
\begin{center}\vspace{1pt}
\setlength{\tabcolsep}{10pt}\renewcommand{\arraystretch}{1.25}\small
\begin{tabular}{c|c}
\textbf{Horas $\downarrow$} & \textbf{Peças $\downarrow$}\\
\hline
2 & 1\\
\rowcolor{rulebg}$x$ & 3\\
\end{tabular}
\end{center}
\exemp{%
  $\dfrac{2}{x} = \dfrac{1}{3} \;\Rightarrow\; 1\cdot x = 2\cdot3$\\
  $x = \hlm{6\text{ horas}}$%
}

\SH{Exemplo 2 — distância e tempo}
120\,km em 2 horas; quantos km em 5 horas?
\begin{center}\vspace{1pt}
\setlength{\tabcolsep}{10pt}\renewcommand{\arraystretch}{1.25}\small
\begin{tabular}{c|c}
\textbf{km $\downarrow$} & \textbf{Horas $\downarrow$}\\
\hline
120 & 2\\
\rowcolor{rulebg}$x$ & 5\\
\end{tabular}
\end{center}
\exemp{%
  $\dfrac{120}{x} = \dfrac{2}{5} \;\Rightarrow\; 2x = 600$\\
  $x = \hlm{300\text{ km}}$%
}

%─────────────────────────────────────────────────────────────
\TW{Inv}{badgegreen}{Simples — Proporção Inversa}{Grandezas em sentidos opostos}

\regra{\textbf{Inversamente proporcionais:} o aumento de uma provoca
a diminuição da outra. A fração da grandeza inversa é
\textbf{invertida} antes do produto cruzado.}

\SH{Exemplo 1 — dias e funcionários}
2 funcionários em 6 dias; quantos dias com 4? \emph{Mais
funcionários, menos dias → inversa.}
\begin{center}\vspace{1pt}
\setlength{\tabcolsep}{8pt}\renewcommand{\arraystretch}{1.25}\small
\begin{tabular}{c|c}
\textbf{Dias $\downarrow$} & \textbf{Func. $\uparrow$}\\
\hline
6 & 2\\
\rowcolor{rulebg}$x$ & 4\\
\end{tabular}
\end{center}
\exemp{%
  Invertendo a fração de Dias:\\
  $\dfrac{x}{6} = \dfrac{2}{4} \;\Rightarrow\; 4x = 12$\\
  $x = \hlm{3\text{ dias}}$%
}

\SH{Exemplo 2 — canos e tempo}
1 cano enche em 12 horas; e 3 canos iguais?
\begin{center}\vspace{1pt}
\setlength{\tabcolsep}{8pt}\renewcommand{\arraystretch}{1.25}\small
\begin{tabular}{c|c}
\textbf{Horas $\downarrow$} & \textbf{Canos $\uparrow$}\\
\hline
12 & 1\\
\rowcolor{rulebg}$x$ & 3\\
\end{tabular}
\end{center}
\exemp{%
  $\dfrac{x}{12} = \dfrac{1}{3} \;\Rightarrow\; 3x = 12$\\
  $x = \hlm{4\text{ horas}}$%
}

%─────────────────────────────────────────────────────────────
\TW{Comp}{badgepurple}{Regra de Três Composta}{Três ou mais grandezas}

\regra{Multiplica-se a fração principal pelo produto das demais.
\textbf{Diretas} entram com $\frac{v_1}{v_2}$; \textbf{inversas}
com a fração \textbf{invertida} $\frac{v_2}{v_1}$.}

\SH{Passo a passo}
\exemp{%
  \textbf{1.} Identifique a grandeza com $x$ (principal).\\
  \textbf{2.} Para cada outra: se aumenta, a principal aumenta ou
  diminui?\\
  \textbf{3.} Monte $\dfrac{\text{princ.}_1}{x} = \prod$ frações
  ($\frac{v_1}{v_2}$ diretas, $\frac{v_2}{v_1}$ inversas).\\
  \textbf{4.} Resolva com produto cruzado.%
}

\SH{Exemplo — pedreiros, paredes e dias}
3 pedreiros constroem 2 paredes em 4 dias. Quantos dias 6 pedreiros
levam para 6 paredes?
\begin{center}\vspace{1pt}
\setlength{\tabcolsep}{7pt}\renewcommand{\arraystretch}{1.25}\small
\begin{tabular}{c|c|c}
\textbf{Dias $\downarrow$} & \textbf{Pedr. $\uparrow$} & \textbf{Pared. $\downarrow$}\\
\hline
4 & 3 & 2\\
\rowcolor{rulebg}$x$ & 6 & 6\\
\end{tabular}
\end{center}
\exemp{%
  Pedreiros \textbf{inversa} → $\tfrac{6}{3}$; \ Paredes
  \textbf{direta} → $\tfrac{2}{6}$:\\[2pt]
  $\dfrac{4}{x} = \dfrac{6}{3}\cdot\dfrac{2}{6}
   = \dfrac{12}{18} = \dfrac{2}{3}$\\
  $2x = 12 \;\Rightarrow\; x = \hlm{6\text{ dias}}$%
}

%─────────────────────────────────────────────────────────────
\TW{\%}{badgegreen}{Porcentagem}{Razão de 100, conversões e aplicações}

\regra{\textbf{Porcentagem} é uma razão de denominador $100$:}
\formula{$p\% = \dfrac{p}{100}$}

\SH{Conversões: fração ↔ decimal ↔ \%}
\exemp{%
  $25\% = \dfrac{25}{100} = 0{,}25$ \;→\; $\dfrac{1}{4} = \hlm{25\%}$\\[3pt]
  $0{,}08 = \dfrac{8}{100} = \hlm{8\%}$\\[3pt]
  $\dfrac{3}{5} = \dfrac{60}{100} = \hlm{60\%}$%
}

\SH{Porcentagem de um valor}
\formula{$\text{resultado} = V \times \dfrac{p}{100}$}
\exemp{%
  $30\%$ de R\$\,250: $\;250\times0{,}30 = \hlm{75}$\\[3pt]
  Que \% é 18 de 120? $\;\dfrac{18}{120}\times100 = \hlm{15\%}$%
}

\SH{Aumento e desconto}
\exemp{%
  \textbf{Aumento} de $p\%$: fator $(1+p/100)$\\
  R\$\,180 com $+15\%$: $\;180\times1{,}15 = \hlm{207}$\\[3pt]
  \textbf{Desconto} de $p\%$: fator $(1-p/100)$\\
  R\$\,350 com $-20\%$: $\;350\times0{,}80 = \hlm{280}$\\[3pt]
  \textbf{Variação:} R\$\,40 → R\$\,52\\
  $\dfrac{52-40}{40}\times100 = \hlm{30\%}$ de aumento%
}

\nota{\textbf{Sucessivos não se somam:} $+20\%$ seguido de $-20\%$
\emph{não} volta ao original — $1{,}20\times0{,}80 = 0{,}96$, ou seja,
$4\%$ abaixo.}

\end{multicols}

\end{document}
